
\documentclass{article}
\usepackage{colortbl}
\usepackage{makecell}
\usepackage{multirow}
\usepackage{supertabular}

\begin{document}

\newcounter{utterance}

\centering \large Interaction Transcript for game `cladder', experiment `full\_v1.5\_default', episode 8207 with Qwen3.6{-}27B.
\vspace{24pt}

{ \footnotesize  \setcounter{utterance}{1}
\setlength{\tabcolsep}{0pt}
\begin{supertabular}{c@{$\;$}|p{.15\linewidth}@{}p{.15\linewidth}p{.15\linewidth}p{.15\linewidth}p{.15\linewidth}p{.15\linewidth}}
    \# & \multicolumn{2}{c}{Player} && \multicolumn{2}{c}{Game Master} \\
    \hline

    \theutterance \stepcounter{utterance}  
    & & & \multicolumn{4}{p{0.6\linewidth}}{
        \cellcolor[rgb]{0.9,0.9,0.9}{
            \makecell[{{p{\linewidth}}}]{
                \texttt{\tiny{[P1$\langle$GM]}}
                \texttt{You are an expert in causal inference. The following question is not a typical commonsense query, but rather a meticulously designed question created by a professor specializing in causal inference, intended to assess the students' mastery of the course content.} \\
\texttt{Imagine a self{-}contained, hypothetical world with only the following conditions, and without any unmentioned factors or causal relationships: Tijv has a direct effect on tijw and xevo. Xevo has a direct effect on gyzp. Tijw has a direct effect on gyzp. For those who are not tijv and are not xevo, the probability of gyzp is 29\%. For those who are not tijv and are xevo, the probability of gyzp is 14\%. For those who are tijv and are not xevo, the probability of gyzp is 11\%. For those who are tijv and are xevo, the probability of gyzp is 16\%. The overall probability of tijv is 88\%. Will xevo increase the chance of gyzp?} \\
\texttt{Start your answer with 'Yes' or 'No', followed by additional reasoning or evidence to support your explanation.} \\
            }
        }
    }
     \\ \\

    \theutterance \stepcounter{utterance}  
    & \multicolumn{4}{p{0.6\linewidth}}{
        \cellcolor[rgb]{0.561,0.737,0.561}{
            \makecell[{{p{\linewidth}}}]{
                \texttt{\tiny{[P1$\rangle$GM]}}
                \texttt{No} \\
\\ 
\texttt{To determine whether `xevo` increases the chance of `gyzp`, we must calculate the Average Causal Effect (ACE) of `xevo` on `gyzp`. In causal inference, this requires comparing the expected outcome of `gyzp` if everyone were set to `xevo=1` versus if everyone were set to `xevo=0`, while accounting for the causal structure.} \\
\\ 
\texttt{**1. Identify the Causal Structure and Confounding**} \\
\texttt{The problem states:} \\
\texttt{*   `Tijv` has a direct effect on `xevo`.} \\
\texttt{*   `Tijv` has a direct effect on `gyzp`.} \\
\texttt{*   `Xevo` has a direct effect on `gyzp`.} \\
\texttt{*   `Tijw` has a direct effect on `gyzp` (but `Tijw` is not mentioned as a cause of `xevo` or affected by `xevo` in a way that creates a backdoor path between `xevo` and `gyzp` other than through `Tijv`? Let's check).} \\
\\ 
\texttt{Actually, let's look at the paths between `Xevo` and `Gyzp`:} \\
\texttt{1.  **Direct Path:** `Xevo {-}> Gyzp`} \\
\texttt{2.  **Backdoor Path:** `Xevo <{-} Tijv {-}> Gyzp`} \\
\\ 
\texttt{Since `Tijv` is a common cause of both `Xevo` and `} \\
            }
        }
    }
    & & \\ \\

    \theutterance \stepcounter{utterance}  
    & & & \multicolumn{2}{p{0.3\linewidth}}{
        \cellcolor[rgb]{0.9,0.9,0.9}{
            \makecell[{{p{\linewidth}}}]{
                \texttt{\tiny{[GM$|$GM]}}
                \texttt{No} \\
            }
        }
    }
    & & \\ \\

    \theutterance \stepcounter{utterance}  
    & & & \multicolumn{2}{p{0.3\linewidth}}{
        \cellcolor[rgb]{0.9,0.9,0.9}{
            \makecell[{{p{\linewidth}}}]{
                \texttt{\tiny{[GM$|$GM]}}
                \texttt{yes} \\
            }
        }
    }
    & & \\ \\

    \theutterance \stepcounter{utterance}  
    & & & \multicolumn{2}{p{0.3\linewidth}}{
        \cellcolor[rgb]{0.9,0.9,0.9}{
            \makecell[{{p{\linewidth}}}]{
                \texttt{\tiny{[GM$|$GM]}}
                \texttt{game\_result = LOSE} \\
            }
        }
    }
    & & \\ \\

\end{supertabular}
}

\end{document}
